\section{Results and Discussion}
\label{sec:results-discussion}

\subsection{Integrated answers to Q1--Q3}

Q1 produced a composition-only surrogate with MAE $2.4715$, Spearman
$0.9383$, and 20 true Top-50 hits on the selection-involved file~03
assessment.  Its low-tail MAE remained $8.2654$, so the model is used for
prioritization rather than physical certification.  For the frozen Q2/Q3
proxy, exhaustive enumeration and the adopted 104-variable matrices selected
Y--Ti--$\mathrm{NO}_2$ ($\widehat V_{23}=1.491928$) and
Na--Nb--$\mathrm O_3$ ($0.540023$), with runner-up margins $0.313819$ and
$0.013123$.  These unlabeled-pool values are numerical proxy rankings, not
statistically resolved material winners.

The Q1 mixture improved both global error and low-tail retrieval relative to
direct HGB: Top-25/Top-50 hits were 8/20 rather than 4/13.  A post-freeze
expected-penalty control had slightly lower file~03 MAE ($2.4624$), but the
preselected mixture performed better in the grouped development folds.  We
therefore retain the frozen model and avoid retrospective replacement from a
single metric.

\subsection{Finite-budget value and solver sampling}

At budgets 10, 25, 50, and 100, the Q1 ranking retrieved 6, 13, 23, and 34
records with observed $V<5$, versus uniform expectations 0.1764, 0.4409,
0.8818, and 1.7637 (Table~\ref{tab:q1-budget}).  This demonstrates value on
the selection-involved replay, not a new-domain success probability.

Across the four adopted matrices, Kaiwu SDK~1.4.3 CPU runs returned 99/1,200
legal samples and hit the group optimum in 99/120 seed-level attempts, with
83.413~s recorded solver time.  Local SA returned 907/1,200 legal samples and
hit in all 120 ten-read batches; audited O15-R01--R04 calls totaled
178.131869~s.  Different implementations, attempt definitions, and timing
boundaries preclude a speed comparison.  The earlier 162.556~s total remains
\textbf{UNVERIFIED}.

\subsection{Four-layer error audit}

The four layers are reported separately because their units and conditioning
differ.  Surrogate error is measured only by the selection-involved Q1
assessment; Q2/Q3 prediction error remains NA.  Mapping controls gave
$(\epsilon_{\mathrm{coef}},\epsilon_{\mathrm{energy}})$ of
$(2.13\times10^{-14},1.42\times10^{-12})$ for 104-variable Q2 and
$(9.95\times10^{-14},4.57\times10^{-12})$ for 208-variable Q3.  Adopted
integer matrices instead had zero formula-to-matrix residual over all 2,704
legal states per group, while their comparable Kaiwu~1.4.4 floating residual
remains NA.

Uniform eight-bit controls retained the raw minimizer only within a tied
lowest tier: six Q2 candidates and 21 or 24 Q3 candidates, depending on the
rounding rule.  Candidate-ID tie-breaking caused score losses $0.6474$ and
$0.6451$.  The adopted two-level mapping keeps a known group argmin with
$\max|h|=102$, $\max|J|=1$, and ground-state gap 2, but is oracle-assisted
and discards non-winner ordering.

The blind-pool models do not receive the optimum identity but use each full
within-pool proxy-score range for scaling.  All eight local-SA runs returned
300/300 legal samples in 38.533--48.536~s; floating/eight-bit hit ranges were
25--27/30 and 11--16/30 batches, with 11--34 minimum-tier ties.  Q2 selected
Y--Ti--$\mathrm{NO}_2$, while Q3 selected Ca--Sn--$\mathrm O_3$ after
original-scale comparison.  These are local-SA proxy results.  Separately,
the four final SPQC tasks had legal fractions 10\%, 40\%, 0\%, and 20\%; three
returned a group proxy argmin.  The small unequal returns are not a stability
estimate or evidence of quantum advantage.

\subsection{Pipeline robustness and measured cost boundary}

The intake audit rejected 13 injected defect classes, and R1--R4 covered
illegal-only output, missing SDK licence, complete Q2/Q3 domains, and source-
field validation.  The complete build--convert--adapt--check pipeline took
6.551~ms for Q2 and 7.430--7.591~ms per Q3 group; the measured three-group Q3
total was 23.197~ms.  Full-workflow SDK times were 20.81--20.96~s per group.
A reported 0.48--0.69~s scope covers only 30 \texttt{opt.solve()} calls and
lacks its raw timing package, so it remains \textbf{UNVERIFIED}.  SPQC
creation-to-finish intervals were 16--112~s; queue, complete execution, and
monetary cost remain NA.  Appendix~A.5 preserves the detailed records.

\subsection{Stability, contribution, and limitations}

Alternative Q1 group allocations retained a mixture-screening advantage
($40.2\%$ and $42.6\%$ versus $28.0\%$ and $26.0\%$ for direct HGB), but
file~03 informed model selection and blind pools remain unlabeled.  The
107-variable projection changed the optimum, integer blind-pool tiers tied,
and four hardware tasks cannot characterize stability.

The contribution is an auditable separation of surrogate prediction,
oracle-assisted full-pool representation, qualified blind-pool search,
mapping precision, quantization, and sampling evidence
\cite{Ward2016Magpie,Chen2024PerovskiteML,Lucas2014IsingNP,Glover2019QUBOTutorial}.
It is not a new universal surrogate or annealing algorithm.  Independent
labels, non-oracle multilevel encodings, reconciled timing and quota records,
the missing probe receipt, and repeated hardware batches remain future work.
